% GENERATED FILE. Edit canonical lessons, then run npm run generate:books.
% canonical-source-hash: fnv1a64:ec7f4017c2dae8c7
% canonical-lessons: FR-C480-fracture, FR-C480-platre, FR-C480-bequille, FR-C480-boiter, FR-C480-os

\chapter{Une jambe dans le plâtre}
\label{ch:fr-une-jambe-dans-le-platre}
\hlchaptermodality{480}

\begin{chapteropening}
\textbf{By the end of this chapter:} \emph{I can talk about a broken bone and getting about.}

Complete the last lesson of chapter 480: I can talk about a broken bone and getting about.
\end{chapteropening}

\section[la fracture]{la fracture — fracture, break (of a bone)}

\label{lesson:FR-C480-fracture}



\begin{quote}
Before the new one: say the French for burn, then the French for bump.
\end{quote}

\subsection*{You'll want to know: la fracture}
\textbf{la fracture} — \textquotedblleft{}fracture, break (of a bone)\textquotedblright{}.

\begin{quote}
La radio montre une fracture du poignet.
\end{quote}

\begin{quote}
\emph{The X-ray shows a broken wrist.}
\end{quote}

The everyday verb is \emph{se casser}: \emph{il s'est cassé la jambe}.

\begin{cousinweb}
\emph{la fracture} is a learned borrowing from Latin \emph{frāctūra} \textquotedblleft{}a break\textquotedblright{}, from \emph{frangere} \textquotedblleft{}to break\textquotedblright{} (Lewis \& Short).
\end{cousinweb}

\begin{grammarlens}[title={What you've built}]
One more word for telling, explaining and describing at B1.
\end{grammarlens}

\subsection*{Your turn}
\begin{itemize}
\raggedright
  \item \emph{Say it:} \emph{la fracture}
  \item \emph{Say it:} \emph{la fracture}, once more
  \item \emph{Say it:} a sentence of your own with \emph{la fracture}
  \item \emph{Recall:} say the French for burn, then the French for bump, then say \emph{la fracture} again
\end{itemize}

\begin{tcolorbox}[breakable,colback=teal!4,colframe=teal!35!black,title={Before you move on}]
What does la fracture mean? (\textquotedblleft{}Fracture, break (of a bone)\textquotedblright{}.) Say it once more.
\end{tcolorbox}
\section[le plâtre]{le plâtre — plaster cast; plaster}

\label{lesson:FR-C480-platre}



\begin{quote}
Before the new one: say the French for bump, then the French for fracture.
\end{quote}

\subsection*{You'll want to know: le plâtre}
\textbf{le plâtre} — \textquotedblleft{}plaster cast; plaster\textquotedblright{}.

\begin{quote}
Mon fils a eu la jambe dans le plâtre pendant six semaines.
\end{quote}

\begin{quote}
\emph{My son had his leg in a cast for six weeks.}
\end{quote}

\emph{Être dans le plâtre} — to be in a cast. Also building plaster.

\begin{cousinweb}
\emph{le plâtre} comes from Old French \emph{plastre}, which goes back to Latin \emph{emplastrum} \textquotedblleft{}a plaster for a wound\textquotedblright{}, a borrowing from Greek \emph{émplastron}.
\end{cousinweb}

\begin{grammarlens}[title={What you've built}]
One more word for telling, explaining and describing at B1.
\end{grammarlens}

\subsection*{Your turn}
\begin{itemize}
\raggedright
  \item \emph{Say it:} \emph{le plâtre}
  \item \emph{Say it:} \emph{le plâtre}, once more
  \item \emph{Say it:} a sentence of your own with \emph{le plâtre}
  \item \emph{Recall:} say the French for bump, then the French for fracture, then say \emph{le plâtre} again
\end{itemize}

\begin{tcolorbox}[breakable,colback=teal!4,colframe=teal!35!black,title={Before you move on}]
What does le plâtre mean? (\textquotedblleft{}Plaster cast; plaster\textquotedblright{}.) Say it once more.
\end{tcolorbox}
\section[la béquille]{la béquille — crutch}

\label{lesson:FR-C480-bequille}



\begin{quote}
Before the new one: say the French for fracture, then the French for plaster cast.
\end{quote}

\subsection*{You'll want to know: la béquille}
\textbf{la béquille} — \textquotedblleft{}crutch\textquotedblright{}.

\begin{quote}
Elle marche avec des béquilles depuis son accident.
\end{quote}

\begin{quote}
\emph{She has been walking on crutches since her accident.}
\end{quote}

Also the kickstand of a bike or scooter.

\begin{cousinweb}
\emph{la béquille} is usually traced to \emph{bec} \textquotedblleft{}beak\textquotedblright{}, for the curved handle; the origin is not fully settled.
\end{cousinweb}

\begin{grammarlens}[title={What you've built}]
One more word for telling, explaining and describing at B1.
\end{grammarlens}

\subsection*{Your turn}
\begin{itemize}
\raggedright
  \item \emph{Say it:} \emph{la béquille}
  \item \emph{Say it:} \emph{la béquille}, once more
  \item \emph{Say it:} a sentence of your own with \emph{la béquille}
  \item \emph{Recall:} say the French for fracture, then the French for plaster cast, then say \emph{la béquille} again
\end{itemize}

\begin{tcolorbox}[breakable,colback=teal!4,colframe=teal!35!black,title={Before you move on}]
What does la béquille mean? (\textquotedblleft{}Crutch\textquotedblright{}.) Say it once more.
\end{tcolorbox}
\section[boiter]{boiter — to limp}

\label{lesson:FR-C480-boiter}



\begin{quote}
Before the new one: say the French for plaster cast, then the French for crutch.
\end{quote}

\subsection*{You'll want to know: boiter}
\textbf{boiter} — \textquotedblleft{}to limp\textquotedblright{}.

\begin{quote}
Depuis sa chute, mon chien boite un peu.
\end{quote}

\begin{quote}
\emph{Since his fall, my dog limps a little.}
\end{quote}

The adjective \emph{boiteux} means limping, and also shaky: \emph{une excuse boiteuse}.

\begin{cousinweb}
The origin of \emph{boiter} is not settled. A link with \emph{boîte}, once the socket of a joint, has been proposed.
\end{cousinweb}

\begin{grammarlens}[title={What you've built}]
One more word for telling, explaining and describing at B1.
\end{grammarlens}

\subsection*{Your turn}
\begin{itemize}
\raggedright
  \item \emph{Say it:} \emph{boiter}
  \item \emph{Say it:} \emph{boiter}, once more
  \item \emph{Say it:} a sentence of your own with \emph{boiter}
  \item \emph{Recall:} say the French for plaster cast, then the French for crutch, then say \emph{boiter} again
\end{itemize}

\begin{tcolorbox}[breakable,colback=teal!4,colframe=teal!35!black,title={Before you move on}]
What does boiter mean? (\textquotedblleft{}To limp\textquotedblright{}.) Say it once more.
\end{tcolorbox}
\section[l'os]{l'os — bone}

\label{lesson:FR-C480-os}



\begin{quote}
Before the new one: say the French for crutch, then the French for to limp.
\end{quote}

\subsection*{You'll want to know: l'os}
\textbf{l'os} — \textquotedblleft{}bone\textquotedblright{}.

\begin{quote}
Le chien ronge un os dans le jardin.
\end{quote}

\begin{quote}
\emph{The dog is gnawing a bone in the garden.}
\end{quote}

Masculine. Singular \emph{un os} with a sounded \emph{s}; plural \emph{des os}, where the \emph{s} is usually silent.

\begin{cousinweb}
\emph{l'os} is inherited from Latin \emph{os}, \emph{ossis} \textquotedblleft{}bone\textquotedblright{} (EtymDB; Lewis \& Short).
\end{cousinweb}

\begin{grammarlens}[title={What you've built}]
One more word for telling, explaining and describing at B1.
\end{grammarlens}

\subsection*{Your turn}
\begin{itemize}
\raggedright
  \item \emph{Say it:} \emph{l'os}
  \item \emph{Say it:} \emph{l'os}, once more
  \item \emph{Say it:} a sentence of your own with \emph{l'os}
  \item \emph{Recall:} say the French for crutch, then the French for to limp, then say \emph{l'os} again
\end{itemize}

\begin{tcolorbox}[breakable,colback=teal!4,colframe=teal!35!black,title={Before you move on}]
What does l'os mean? (\textquotedblleft{}Bone\textquotedblright{}.) Say it once more.
\end{tcolorbox}
